% Notas de aula — Capítulo 3: Termodinâmica da Atmosfera
% Curso: Meteorologia Prática (baseado em R. Stull, Practical Meteorology, CC BY-NC-SA 4.0)
% Caixas verdes = conteúdo literal do quadro; linhas verdes = encenação.
\documentclass[11pt]{article}
\usepackage{../comum/preambulo}
\graphicspath{{../figuras/cap03/}}

\title{Capítulo 3 --- Termodinâmica da Atmosfera\\
{\large Notas de aula (roteiro de quadro-negro)}}
\author{Curso de Meteorologia Prática}
\date{}

\begin{document}
\maketitle
\thispagestyle{fancy}

\noindent\textbf{Plano sugerido:} 2 aulas de 2\,h.
\emph{Aula 1:} \S\ref{sec:primeira}--\S\ref{sec:theta} (1ª lei, processos
adiabáticos, temperatura potencial).
\emph{Aula 2:} \S\ref{sec:euler}--\S\ref{sec:superficie} (balanço euleriano
de calor, balanço de superfície, razão de Bowen).

\section{Calor e energia interna}

\begin{noquadro}[abertura]
\textbf{Cap.~3 --- Termodinâmica}\\[2pt]
``Calor não é substância; é energia em trânsito.''\\[2pt]
O ar armazena energia interna ($\propto T$) e troca por três vias:\\
radiação \quad $|$ \quad condução/turbulência \quad $|$ \quad calor latente
\end{noquadro}

A via radiativa acabamos de quantificar \javimos{2}{balanço radiativo};
a turbulenta será o coração da camada limite \veremos{18}{fluxos
turbulentos}; a latente depende da água e domina os trópicos
\veremos{4}{calor latente} \veremos{7}{precipitação}. Hoje montamos a
contabilidade que une as três: a primeira lei aplicada ao ar.

\begin{formalivro}[calores específicos]
\[ C_p = 1004\ \mathrm{J\,kg^{-1}K^{-1}}\ (\text{pressão const.}),\quad
   C_v = 717\ \mathrm{J\,kg^{-1}K^{-1}}\ (\text{volume const.}),\quad
   C_p = C_v + \Re_d. \]
\end{formalivro}

Por que $C_p > C_v$? A pressão constante, parte do calor vira trabalho de
expansão — é a mesma diferença $\Re_d$ da lei dos gases
\javimos{1}{lei dos gases}. Na atmosfera quase tudo acontece ``a pressão
imposta pelo ambiente'', então $C_p$ é o calor específico do dia a dia.

\begin{noquadro}[a água como bateria térmica]
$L_v = 2{,}5\times10^{6}$\,J/kg\\[2pt]
evaporar 1\,g de água ``rouba'' o calor que esquentaria
\textbf{2{,}5\,kg de ar} em 1\,K.
\end{noquadro}

Esse número gigante é o motivo de o suor funcionar, de a maritimidade
amaciar climas — e de uma tempestade tropical ser uma máquina movida a
vapor \veremos{16}{ciclones tropicais}.

\section{Primeira lei da termodinâmica}\label{sec:primeira}

\quadro{Deduzir a partir de $dU = \delta q\,m - P\,dV$; chegar na forma
por unidade de massa do livro. Desenhar o pistão (Fig.~3.2 do livro): calor
entra, parte esquenta o gás, parte empurra o pistão.}

\begin{formadif}[forma diferencial]
Por unidade de massa, com volume específico $\alpha = 1/\rho$:
\[ \delta q = c_v\,dT + P\,d\alpha
 \;\;\overset{P\alpha = \Re_d T}{=}\;\; c_p\,dT - \alpha\,dP. \]
A segunda forma é a útil em meteorologia: pressão — não volume — é a
coordenada natural de quem sobe e desce na atmosfera
\javimos{1}{perfil de pressão}.
\end{formadif}

\begin{formalivro}[Stull eq.~3.2d]
\[ \Delta q = C_p\,\Delta T - \frac{1}{\rho}\,\Delta P. \]
Aquecer o ar ($\Delta q$) vai parte para esquentar ($C_p\Delta T$), parte
para trabalho de expansão ($-\Delta P/\rho$).
\end{formalivro}

\section{Processos adiabáticos e gradiente adiabático seco}\label{sec:adiab}

\quadro{Pergunta à turma ANTES de deduzir: por que uma parcela que sobe
esfria, se ninguém tirou calor dela? (Deixar responderem; a resposta certa
é ``ela paga o trabalho de expansão com energia interna''.)}

Uma parcela que sobe encontra pressão menor \javimos{1}{$P(z)$
exponencial}, expande e faz trabalho contra o ambiente — esfria \emph{sem
trocar calor}. O ar é péssimo condutor: subidas de minutos a horas são
adiabáticas em excelente aproximação \javimos{1}{terminologia de
processos}.

\begin{formadif}[dedução do gradiente adiabático seco]
Adiabático: $\delta q = 0 \Rightarrow c_p\,dT = \alpha\,dP$. Com a
hidrostática $dP = -\rho g\,dz$ \javimos{1}{equilíbrio hidrostático}:
\[ c_p\,dT = -g\,dz
\quad\Longrightarrow\quad
\boxed{\ \frac{dT}{dz}\bigg|_{ad} = -\frac{g}{c_p} \equiv -\Gamma_d\ } \]
\end{formadif}

\begin{formalivro}[Stull eq.~3.14]
\[ \Gamma_d = \frac{g}{C_p} = 9{,}8\ \mathrm{K\,km^{-1}}
\qquad \text{(gradiente adiabático seco)}. \]
\end{formalivro}

\begin{noquadro}[não confundir os dois gradientes]
$\Gamma_d = 9{,}8$\,K/km: o que a \emph{parcela} faz ao subir
(física, fixo)\\
$\gamma \simeq 6{,}5$\,K/km: o que a \emph{sondagem} mostra
(ambiente, varia dia a dia)\\[2pt]
comparar os dois $=$ estabilidade $\to$ Cap.~5
\end{noquadro}

Essa distinção parcela/ambiente é o erro conceitual mais comum do curso —
vale dez minutos de aula. A comparação decide se o ar afunda de volta
(estável) ou dispara para cima (instável) \veremos{5}{estabilidade
estática}; quando a parcela satura, o $\Gamma$ dela muda
\veremos{4}{saturação} \veremos{5}{gradiente adiabático úmido}.

\section{Temperatura potencial}\label{sec:theta}

\quadro{Motivar: ``queremos um termômetro que não mude só porque a parcela
subiu''. Deduzir $\theta$ no quadro a partir da mesma linha adiabática.}

\begin{formadif}[dedução]
De $c_p\,dT = \alpha\,dP = \Re_d T\,dP/P$:
\[ \frac{dT}{T} = \frac{\Re_d}{c_p}\frac{dP}{P}
\;\Longrightarrow\;
T\,P^{-\Re_d/c_p} = \text{const ao longo da adiabática}. \]
Definindo a constante pela pressão de referência $P_0 = 100$\,kPa:
\end{formadif}

\begin{formalivro}[Stull eq.~3.12]
\[ \boxed{\ \theta = T\left(\frac{P_0}{P}\right)^{\Re_d/C_p}\ },
\qquad \frac{\Re_d}{C_p} = 0{,}28571. \]
$\theta$ = temperatura que a parcela teria se trazida adiabaticamente a
100\,kPa. \textbf{Conservada} em processos adiabáticos secos.
\end{formalivro}

\begin{noquadro}[por que $\theta$ vale ouro]
$T$ muda só de subir; $\theta$ não: é a \emph{etiqueta} da parcela.\\
comparar ar de alturas diferentes $\Rightarrow$ comparar $\theta$, nunca
$T$.\\
aproximação de bolso: $\theta \simeq T + (9{,}8\ \mathrm{K/km})\,z$.
\end{noquadro}

\begin{exemplo}[ar de altitude que desce a serra]
Ar a $P = 70$\,kPa e $T = 5\,°$C (278\,K) desce adiabaticamente até
$P = 100$\,kPa (litoral):
\[ \theta = 278\times(100/70)^{0{,}2857} = 278\times1{,}107 = 308\ \mathrm{K}
 = 35\,°\mathrm{C}. \]
O mesmo ar que estava a $5\,°$C chega \emph{quente e seco} ao nível do
mar — mecanismo dos ventos tipo föhn \veremos{17}{ventos de montanha}.
\emph{Verificação: $\Gamma_d\Delta z \simeq 9{,}8\times3 \simeq 29$\,K
$\Rightarrow 278+29=307$\,K, consistente.}
\end{exemplo}

Conexão com a física: $\theta$ é a entropia disfarçada —
$ds = c_p\,d\ln\theta$ (exercício T3). Superfícies isentrópicas
($\theta$ const) são as trilhas naturais do escoamento: o ar viaja por
elas, e as cartas isentrópicas exploram isso \veremos{12}{análise
isentrópica}. Nos diagramas termodinâmicos, as adiabáticas secas são as
linhas de $\theta$ constante \veremos{5}{diagramas termodinâmicos}. Para
ar úmido usa-se $\theta_v$ (com a correção virtual \javimos{1}{temperatura
virtual}).

\section{Balanço euleriano de calor}\label{sec:euler}

\quadro{Desenhar um cubo FIXO no espaço com vento atravessando; setas de
entrada/saída de calor. Contraste com a parcela que viaja (Lagrange):
duas janelas para a mesma física.}

Num ponto fixo (visão \emph{euleriana} — a da estação meteorológica e a
do modelo numérico), a temperatura muda por advecção, radiação,
turbulência e calor latente:

\begin{formadif}[derivada material]
A ponte entre parcela (Lagrange) e ponto fixo (Euler) é a regra da cadeia:
\[ \frac{DT}{Dt} = \frac{\partial T}{\partial t}
 + u\frac{\partial T}{\partial x} + v\frac{\partial T}{\partial y}
 + w\frac{\partial T}{\partial z}
\;\Longrightarrow\;
\frac{\partial T}{\partial t} = -\,\vec{U}\cdot\nabla T + \text{fontes}. \]
O termo $-\vec U\cdot\nabla T$ é a \textbf{advecção}: vento soprando de
região quente esquenta o ponto. Esta mesma estrutura reaparecerá para
umidade \veremos{4}{balanço de umidade}, momento \veremos{10}{equações do
movimento} e qualquer traçador \veremos{19}{dispersão de poluentes}.
\end{formadif}

\begin{formalivro}[Stull eq.~3.51, simplificada]
\[ \frac{\Delta T}{\Delta t} \simeq
 \underbrace{-\,U\frac{\Delta T}{\Delta x}}_{\text{advecção}}
 \;-\; \underbrace{\frac{1}{\rho C_p}\frac{\Delta F_z}{\Delta z}}_{\substack{\text{divergência de fluxo}\\ \text{(rad.\ + turb.)}}}
 \;+\; \underbrace{\frac{L_v}{C_p}\frac{\Delta r_{cond}}{\Delta t}}_{\text{calor latente}}. \]
\end{formalivro}

\begin{noquadro}[exemplo de advecção --- fazer a conta]
vento $U = 10$\,m/s, gradiente $5$\,K$/100$\,km:\\[2pt]
$\partial T/\partial t = -U\,\partial T/\partial x
 = 10\times\dfrac{5}{10^5}$\,K/s $= 5\times10^{-4}$\,K/s
 $\simeq \mathbf{1{,}8}$\,\textbf{K/h}
\end{noquadro}

É a frente fria chegando \veremos{12}{frentes}: quedas de 10\,K em poucas
horas são advecção quase pura. A simulação do Caso~3 do notebook faz essa
frente viajar — e mostra a difusão numérica que os esquemas de advecção
dos modelos precisam domar \veremos{20}{esquemas numéricos}.

\section{Balanço de calor na superfície}\label{sec:superficie}

\quadro{Desenhar a superfície com as 4 setas: $F^*$ (radiação líquida,
já conhecida), $F_H$ (sensível), $F_E$ (latente), $F_G$ (solo). Perguntar:
para onde vai a energia do Sol ao meio-dia?}

\begin{formalivro}[Stull eq.~3.52]
\[ F^{*} = F_H + F_E + F_G, \]
com a \textbf{razão de Bowen} $\beta = F_H/F_E$: deserto $\beta\sim5$;
gramado $\sim0{,}5$; floresta tropical $\sim0{,}2$; oceano $\sim0{,}1$.
\end{formalivro}

\begin{noquadro}[mesma entrada, climas opostos]
$F^*$ igual, partições diferentes:\\
deserto: $F^* \to F_H$ \ (ar seco e MUITO quente: tardes de 45\,°C)\\
floresta: $F^* \to F_E$ \ (evapora: ameno e úmido; a chuva de Manaus)\\[2pt]
razão de Bowen $\beta = F_H/F_E$ é o ``caráter'' da superfície.
\end{noquadro}

O $F^*$ vem do capítulo anterior \javimos{2}{radiação líquida}; a partição
em $F_H$ vs.\ $F_E$ decide a profundidade da camada limite convectiva
\veremos{18}{camada limite convectiva} e o combustível disponível para
tempestades ($F_E$ carrega vapor que vira CAPE)
\veremos{14}{tempestades}. Simulação do ciclo diurno no Caso~2 do
notebook; o exercício T5 compara deserto e floresta.

Fluxo sensível pela fórmula de transferência (\emph{bulk}):
$F_H = \rho C_p\,C_H\,M\,(\theta_{sfc} - \theta_{ar})$, com
$C_H \sim 2\times10^{-3}$ e $M$ o vento — note que já é $\theta$, não $T$,
que entra \javimos{3}{temperatura potencial, \S\ref{sec:theta}}. A mesma
forma \emph{bulk} dará a evaporação \veremos{4}{fórmula bulk da
evaporação} e o atrito \veremos{18}{arrasto na superfície}.

\section{Temperaturas aparentes e sensores}

O corpo humano não mede $T$: mede \emph{fluxo de calor} — e vento e
umidade o alteram (Stull \S3.7).

\begin{formalivro}[índice de resfriamento pelo vento (wind chill)]
\[ T_{wc} = 13{,}12 + 0{,}6215\,T - 11{,}37\,M^{0{,}16}
 + 0{,}3965\,T\,M^{0{,}16}, \]
com $T$ em °C e $M$ em km/h (vento a 10\,m; fórmula JAG/TI usada
pelo Stull). O vento aumenta o coeficiente de troca da fórmula
\emph{bulk} \javimos{3}{$F_H$ bulk} e ``rouba'' calor mais rápido —
$0\,°$C com 50\,km/h sente-se como $-8\,°$C.
\end{formalivro}

No calor, o análogo é o índice de calor/humidex: umidade alta trava a
evaporação do suor — o $T_w$ do próximo capítulo é o limite físico
\veremos{4}{bulbo úmido}. \textbf{Sensores de $T$} (Stull \S3.8):
termômetros de líquido, termopares, termistores e resistências de
platina (padrão das estações automáticas) — todos exigem abrigo
ventilado (o erro clássico é medir radiação, não temperatura
\javimos{2}{balanço radiativo do sensor}).

\section{Síntese da aula}

\begin{noquadro}[mapa final --- canto do quadro]
\[ \delta q = c_p dT - \alpha dP \;\to\; \Gamma_d = \frac{g}{c_p}
 \;\to\; \theta = T\Big(\frac{P_0}{P}\Big)^{0{,}286} \]
\[ \frac{\partial T}{\partial t} = -\vec U\cdot\nabla T + \cdots
 \qquad F^* = F_H + F_E + F_G \]
\end{noquadro}

A 1ª lei vira: régua de parcela ($\Gamma_d$), etiqueta conservada
($\theta$), equação de previsão (euleriana — o embrião do Cap.~20) e
contabilidade de superfície (Bowen). O próximo capítulo acrescenta o
ingrediente que falta — a água e seu calor latente
\veremos{4}{vapor d'água}.

\section{Exercícios complementares}\label{sec:exercicios}

Complementam a seção de exercícios do livro. Soluções (professor) em
\texttt{cap03\_solucoes.ipynb}.

\subsection*{Teóricos}

\begin{enumerate}[label=\textbf{T\arabic*.},itemsep=4pt]
  \item Deduza $\Gamma_d = g/c_p$ a partir da 1ª lei e da hidrostática, e
        explique por que $\Gamma_d$ não depende da altura (nesta
        aproximação). Calcule $\Gamma_d$ em Marte
        ($g = 3{,}7$\,m/s$^2$, $c_p^{CO_2} = 845$\,J\,kg$^{-1}$K$^{-1}$).
  \item Integre $c_p\,dT/T = \Re_d\,dP/P$ e obtenha
        $\theta = T(P_0/P)^{\Re_d/c_p}$; mostre que
        $\theta \simeq T + \Gamma_d z$ para pequenas excursões verticais.
  \item Mostre que $ds = c_p\,d\ln\theta$, isto é, que superfícies
        isentrópicas são superfícies de $\theta$ constante — e conclua que
        processos adiabáticos reversíveis conservam $\theta$.
  \item Advecção pura: para
        $\partial T/\partial t = -U\,\partial T/\partial x$ com $U$
        constante, verifique que $T(x,t) = f(x-Ut)$ é solução para
        qualquer perfil $f$ — a forma viaja sem se deformar.
  \item Numa tarde com $F^* = 500$\,W/m$^2$ e $F_G = 50$\,W/m$^2$, compare
        $F_H$ e $F_E$ para $\beta = 5$ (deserto) e $\beta = 0{,}2$
        (floresta), e estime quanta água a floresta evapora por hora e
        por m$^2$.
\end{enumerate}

\subsection*{Numéricos (no notebook)}

\begin{enumerate}[label=\textbf{N\arabic*.},itemsep=4pt]
  \item Integre a subida adiabática de uma parcela como EDO
        ($dT/dz = -g/c_p$) e verifique que $\theta$ calculada com a
        fórmula exata permanece constante ao longo da trajetória (compare
        também com $\theta \simeq T+\Gamma_d z$).
  \item No modelo de ciclo diurno do Caso~2, compare uma superfície
        ``continental'' (capacidade térmica baixa) com uma ``oceânica''
        (alta): amplitude e defasagem do máximo de temperatura.
  \item Repita o experimento de advecção do Caso~3 com esquema
        \emph{upwind} em resoluções $\Delta x$ diferentes e quantifique a
        difusão numérica (largura da frente após 24\,h vs.\ resolução).
  \item Calcule $\theta$ e $\theta_v$ da sondagem do Cap.~1 com o MetPy e
        identifique a camada de mistura (onde $\theta \approx$ constante)
        e as camadas estáveis.
\end{enumerate}

\vspace{1em}
\noindent\rule{\linewidth}{0.4pt}\\
{\scriptsize Material derivado de R.~Stull, \emph{Practical Meteorology}
(CC BY-NC-SA 4.0); este documento herda a mesma licença.}

\end{document}
